% ID: SEL-D0027-OA-E04
% ESTADO: cotejada. Segunda pasada visual completa y borrador compilado revisados.
% Fuente: Selectividad/matematicas_aplicadas/suplemente_septiembre_EXAMEN 6-ANDALUCÍA.pdf
% SHA-256: 601fd1292e7d6b21ceecde2856a283e0e929d02d4e314d97fd7062bda846650a
% Página: 1; opción A; ejercicio original 4.

Se desea estimar la proporción de bares y restaurantes que en el camino de Santiago ofertan el menú del peregrino con un precio máximo de 12 \texteuro{}. Para ello se eligen aleatoriamente 120 establecimientos que ofrecen este menú, de los que 80 tienen un precio máximo de 12 \texteuro{}.
\begin{enumerate}[label=\alph*)]
\item \textbf{(1.6 puntos)} Con un nivel de confianza del 92\%, obtenga el intervalo de confianza para proporción de establecimientos que tienen un precio máximo de 12 \texteuro{}.
\item \textbf{(0.4 puntos)} Si aumentamos el nivel de confianza al 99\%, \textquestiondown{}qué efecto se produce en el error de estimación?
\item \textbf{(0.5 puntos)} \textquestiondown{}Cuántos establecimientos, como mínimo, deberíamos seleccionar para que, con un nivel de confianza del 99\%, el error de la estimación no sea superior a 0.04?
\end{enumerate}
